% Editable resume content; LaTeX is the source of truth.
\resumeProjectLine{\href{https://github.com/ryusudol/orbit-explorer}{\resumeEmphasis{Orbit Explorer}}}{\textit{React, TypeScript, Three.js, D3, Zustand, Vite, WebAssembly}}
\begin{itemize}
  \item \resumeEmphasis{3D Visualization}: Built a browser-based orbital digital twin in \resumeEmphasis{React} and \resumeEmphasis{Three.js} that renders all \resumeEmphasis{32,583} tracked orbital objects at \resumeEmphasis{60~fps} with instancing (WebGPU with a WebGL 2 fallback) and flies to every planet and \resumeEmphasis{450+} moons, with \resumeEmphasis{D3} analysis panels.
  \item \resumeEmphasis{Main-Thread Performance}: Ran orbit propagation in \resumeEmphasis{Rust/WebAssembly} across \resumeEmphasis{4 Web Workers} sharing memory, and kept simulation state out of React, which reads \resumeEmphasis{Zustand} summaries 5 times per second, holding per-frame JavaScript to \resumeEmphasis{2.6 ms} of a 16.7~ms budget.
  \item \resumeEmphasis{Progressive Loading}: Streamed imagery in stages from 2k images to compressed GPU textures and 16k tiles, loading tiles facing the camera first and each world's tiles only on arrival (\textasciitilde80 MB in all), to reach first interaction in \resumeEmphasis{1.2 s} at 50 Mbps.
  \item \resumeEmphasis{Algorithm Design}: Screened 530M object pairs for close approaches \resumeEmphasis{423\ensuremath{\times}} faster than brute force (638 to 1.5 ms) with a spatial grid, and rebuilt a 17,719-link Starlink laser mesh with Dijkstra routing in \resumeEmphasis{3 ms} (London--New York 21~vs.\ 27~ms by ideal fiber).
\end{itemize}
